\chapter{Vertices and edges of the convex hull of a finite set}
\label{chap:prereq_polytope}

Throughout, $\X$ is a finite subset of a real inner product space of finite dimension, $\V$ its set
of vertices and $\I$ its set of adjacent pairs (Definition~\ref{def:adjacency}).

\begin{lemma}[Vertices are points of the set]
  \label{lem:vertices_subset}
  \lean{Set.vertices_subset, Set.Finite.finite_adjacentTo}
  \leanok
  \uses{def:adjacency}
  $\V \subseteq \X$; in particular $\V$ and $\I$ are finite.
\end{lemma}

\begin{proof}
  \leanok
  An exposed point is an extreme point, and the extreme points of $\conv \X$ lie in $\X$.
\end{proof}

\begin{lemma}[Maximizers over a convex hull]
  \label{lem:max_face}
  \lean{Set.mem_convexHull_setOf_inner_eq}
  \leanok
  Let $S$ be any subset of a real inner product space, $w$ a vector and $c \in \R$ with
  $y^\top w \le c$ for every $y \in S$. If $z \in \conv S$ and $z^\top w \ge c$, then $z$ is a
  convex combination of the points $y \in S$ with $y^\top w = c$.
\end{lemma}

\begin{proof}
  \leanok
  Split $S = S_= \cup S_<$ according to whether $y^\top w = c$ or $y^\top w < c$. If $S_<$ is
  empty there is nothing to prove; if $S_=$ is empty, $\conv S$ lies in the convex set
  $\{u : u^\top w < c\}$, a contradiction. Otherwise $\conv S$ is the union of the segments
  $[a, b]$ with $a \in \conv S_=$ and $b \in \conv S_<$; writing $z = (1 - t) a + t b$, we get
  $c \le z^\top w = (1 - t) c + t\, b^\top w$ with $b^\top w < c$, so $t = 0$ and
  $z = a \in \conv S_=$.
\end{proof}

\begin{lemma}[Vertices of a finite set]
  \label{lem:strict_expose}
  \lean{Set.mem_vertices_iff_exists_inner_lt, Set.mem_vertices_of_inner_lt,
    Set.exists_inner_lt_of_mem_vertices, Set.Finite.mem_vertices_of_mem_extremePoints,
    Set.Finite.convexHull_vertices}
  \leanok
  \uses{lem:vertices_subset, lem:max_face}
  For $x \in \X$, the following are equivalent: $x$ is a vertex; there is $w$ with
  $y^\top w < x^\top w$ for every $y \in \X \setminus \{x\}$; $x$ is an extreme point of
  $\conv \X$. Moreover $\conv \X = \conv \V$, and every point of $\X$ is a convex combination of
  vertices.
\end{lemma}

\begin{proof}
  \leanok
  \uses{lem:vertices_subset, lem:max_face}
  If $x$ is exposed by $w$, it is the unique maximizer of $w$ over $\conv \X$, hence strictly better
  than the other points of $\X$. If $x$ strictly beats $\X \setminus \{x\}$ for $w$, every
  maximizer of $w$ over $\conv \X$ is a convex combination of the maximizers in $\X$
  (Lemma~\ref{lem:max_face}), that is of $x$ alone: $x$ is the unique maximizer, an exposed point
  (this direction does not use the finiteness of $\X$). If $x$ is an extreme point of $\conv \X$,
  then $x \in \X$ and $x \notin \conv(\X \setminus \{x\})$ (otherwise $x$ would be an extreme point
  of $\conv(\X \setminus \{x\})$, hence a point of $\X \setminus \{x\}$); separate $x$ strictly from
  this compact convex set (Hahn--Banach) to get $w$. Finally, by the Krein--Milman theorem,
  $\conv \X$ (compact and convex) is the closure of the convex hull of its extreme points, which
  are vertices; as $\V$ is finite, $\conv \V$ is closed, so $\conv \X \subseteq \conv \V$.
\end{proof}

\begin{lemma}[Linear forms are maximized at vertices]
  \label{lem:max_at_vertex}
  \lean{Set.Finite.exists_mem_vertices_forall_inner_le}
  \leanok
  \uses{lem:strict_expose}
  If $\X$ is nonempty, for every $\theta$ there is a vertex $v$ with $y^\top \theta \le v^\top \theta$
  for every $y \in \X$.
\end{lemma}

\begin{proof}
  \leanok
  \uses{lem:strict_expose}
  $\V$ is finite and nonempty ($\conv \V = \conv \X \ne \emptyset$); take $v \in \V$ maximizing
  $\theta$ over $\V$. The half-space $\{u : u^\top \theta \le v^\top \theta\}$ is convex and
  contains $\V$, hence $\conv \V = \conv \X \supseteq \X$.
\end{proof}

\begin{lemma}[Lexicographic maximizers]
  \label{lem:lex_expose}
  \lean{Set.Finite.exists_inner_lt_of_lex}
  \leanok
  Let $x$ maximize $\varphi$ over $\X$, and let $y^\top \psi < x^\top \psi$ for every
  $y \in \X \setminus \{x\}$ with $y^\top \varphi = x^\top \varphi$. Then for some $\delta > 0$,
  $y^\top (\varphi + \delta \psi) < x^\top (\varphi + \delta \psi)$ for every $y \in \X \setminus \{x\}$.
\end{lemma}

\begin{proof}
  \leanok
  For each of the finitely many $y \in \X \setminus \{x\}$, the inequality holds for all small
  enough $\delta > 0$: for every $\delta > 0$ if $y^\top \varphi = x^\top \varphi$, and by
  continuity in $\delta$ if $y^\top \varphi < x^\top \varphi$.
\end{proof}

\begin{lemma}[Adjacency through a direction]
  \label{lem:is_adjacent_iff}
  \lean{Set.Finite.isAdjacent_iff_exists_inner}
  \leanok
  \uses{lem:strict_expose, lem:max_face}
  Two distinct vertices $x, x'$ are adjacent if and only if there is $w$ with
  $x^\top w = x'^\top w$ and $y^\top w < x^\top w$ for every vertex $y \notin \{x, x'\}$. The only
  vertices in $[x, x']$ are $x$ and $x'$. In particular, if $x$ and $x'$ are adjacent there are
  $w$ and $\eps > 0$ with $x^\top w = x'^\top w \ge y^\top w + \eps$ for every vertex
  $y \notin \{x, x'\}$ (the paper's Eq.~(adjacentdef)).
\end{lemma}

\begin{proof}
  \leanok
  \uses{lem:strict_expose, lem:max_face}
  If $[x, x']$ is the set of maximizers over $\conv \X$ of the linear form $w$, then $x$ and $x'$ are
  both maximizers, and a vertex $y \notin \{x, x'\}$ that is a maximizer lies in $[x, x']$, hence in
  the open segment: impossible for an extreme point of $\conv \X$. Conversely, all vertices have
  $y^\top w \le x^\top w$, hence so does $\conv \V = \conv \X$ (Lemma~\ref{lem:strict_expose}), and
  $[x, x']$ consists of maximizers; a maximizer over $\conv \X = \conv \V$ is a convex combination of
  the maximizing vertices (Lemma~\ref{lem:max_face}), which are $x$ and $x'$, so it lies in
  $[x, x']$. For the last claim take $\eps$ the smallest gap over the finitely many vertices.
\end{proof}

\begin{lemma}[Local optimality implies global optimality]
  \label{lem:local_global}
  \lean{Set.Finite.exists_mem_adjacentTo_inner_pos, Set.Finite.adjacentTo_nonempty}
  \leanok
  \uses{lem:strict_expose, lem:max_at_vertex, lem:lex_expose, lem:is_adjacent_iff}
  Let $x \in \V$.
  \begin{enumerate}
    \item If $(y - x)^\top \theta > 0$ for some $y \in \X$, then $(z - x)^\top \theta > 0$ for some
      $z \in \I^x$.
    \item If $\X$ has a point other than $x$, then $\I^x \ne \emptyset$.
  \end{enumerate}
\end{lemma}

\begin{proof}
  \leanok
  \uses{lem:strict_expose, lem:max_at_vertex, lem:lex_expose, lem:is_adjacent_iff}
  Claim 2 follows from claim 1 with $\theta = y - x$. For claim 1, let $w$ strictly expose $x$
  (Lemma~\ref{lem:strict_expose}). Every $y \in \X \setminus \{x\}$ is $y = x + d_y a_y$ with
  $d_y = (x - y)^\top w > 0$ and $a_y = (y - x)/d_y$, so that $a_y^\top w = -1$.

  By Lemma~\ref{lem:max_at_vertex} applied to the finite set $A = \{a_y : y \in \X \setminus \{x\}\}$,
  some vertex $b$ of $A$ maximizes $\theta$ over $A$; in particular
  $b^\top \theta \ge a_{y}^\top \theta = (y - x)^\top \theta / d_y > 0$ for the given $y$. Let $\psi$
  strictly expose $b$ in $A$, and among the points $y'$ with $a_{y'} = b$ (the points of $\X$ on the
  ray $x + \R_+ b$) let $z$ maximize $d_{y'}$. Set $\varphi = \psi + (b^\top \psi)\, w$. For
  $y' \in \X \setminus \{x\}$,
  \[
    (y' - x)^\top \varphi = d_{y'} \bigl(a_{y'}^\top \psi - b^\top \psi\bigr) \le 0,
  \]
  with equality if and only if $a_{y'} = b$. So $x$ and $z$ maximize $\varphi$ over $\X$, and every
  other maximizer $y'$ lies on the ray with $d_{y'} < d_z$, that is in the open segment $(x, z)$.

  $z$ is a vertex: it maximizes $\varphi$, and among the maximizers of $\varphi$ it is the unique
  maximizer of $-w$ ($(-w)^\top y' = (-w)^\top x + d_{y'}$), so Lemma~\ref{lem:lex_expose} gives a
  direction strictly exposing it. $x$ and $z$ are adjacent by Lemma~\ref{lem:is_adjacent_iff} with
  the direction $\varphi$: a vertex $v \notin \{x, z\}$ maximizing $\varphi$ would lie in the open
  segment $(x, z)$, which is impossible for an extreme point of $\conv \X$. Finally
  $(z - x)^\top \theta = d_z\, b^\top \theta > 0$.
\end{proof}

\begin{lemma}[Adjacency lemma, weak form]
  \label{lem:adjacency_ge}
  \lean{Set.Finite.exists_mem_adjacentTo_inner_nonneg}
  \leanok
  \uses{lem:local_global}
  Let $x \in \V$. If $(y - x)^\top \theta \ge 0$ for some $y \in \X \setminus \{x\}$, then
  $(z - x)^\top \theta \ge 0$ for some $z \in \I^x$.
\end{lemma}

\begin{proof}
  \leanok
  \uses{lem:local_global}
  Otherwise $(z - x)^\top \theta < 0$ for every $z$ in the finite set $\I^x$, hence
  $(z - x)^\top (\theta + \eps (y - x)) < 0$ for every $z \in \I^x$ and every small enough
  $\eps > 0$. But $(y - x)^\top(\theta + \eps (y - x)) \ge \eps \norm{y - x}^2 > 0$, and
  Lemma~\ref{lem:local_global} gives $z \in \I^x$ with $(z - x)^\top (\theta + \eps(y - x)) > 0$.
\end{proof}

\begin{lemma}[Cone property]
  \label{lem:cone_form}
  \lean{Set.Finite.smul_mem_convexHull_image_adjacentTo}
  \leanok
  \uses{lem:local_global}
  Let $\theta \in \R^d$ and $x_\star \in \X$ with $y^\top \theta < x_\star^\top \theta$ for every
  $y \in \X \setminus \{x_\star\}$. For $y \in \X \setminus \{x_\star\}$ let
  $u_y = (x_\star - y) / (x_\star - y)^\top \theta$. Then for every $x \in \X \setminus \{x_\star\}$,
  $u_x \in \conv\{u_z : z \in \I^{x_\star}\}$.
\end{lemma}

\begin{proof}
  \leanok
  \uses{lem:local_global, lem:strict_expose}
  $x_\star$ is a vertex (Lemma~\ref{lem:strict_expose}). The set $K = \conv\{u_z : z \in
  \I^{x_\star}\}$ is compact and convex, and lies with $u_x$ in the affine hyperplane
  $\{u : u^\top \theta = 1\}$. If $u_x \notin K$, a strict separation gives $\phi$ and $c$ with
  $\phi^\top u_x > c > \phi^\top u$ for all $u \in K$; then $\psi = c\,\theta - \phi$ satisfies
  $\psi^\top u_z > 0$ for $z \in \I^{x_\star}$ and $\psi^\top u_x < 0$, that is
  $(z - x_\star)^\top \psi < 0$ for every $z \in \I^{x_\star}$ and $(x - x_\star)^\top \psi > 0$,
  contradicting Lemma~\ref{lem:local_global}. (If $\I^{x_\star} = \emptyset$ the same argument
  applies with $K = \emptyset$.)
\end{proof}
